\documentclass[10pt,a4paper]{amsart}
\usepackage[margin=19mm]{geometry}
\usepackage{amsmath,amssymb,mathtools}
\usepackage[scr=rsfs,cal=euler]{mathalfa}
\usepackage{xfrac}
\usepackage[unicode,hidelinks,bookmarks=false]{hyperref}
\newcommand{\ie}{i.e.\ }
\newcommand{\cf}{cf.\ }
\newcommand{\resp}{respectively\ }
\newcommand{\Subsection}{\subsection}
\providecommand{\nobreakdash}{\nobreak\hbox{-}}
\setlength{\emergencystretch}{2em}
\multlinegap0pt
\usepackage{enumitem}
\usepackage{fancyhdr}
\usepackage{mathrsfs}
\usepackage{stmaryrd}
\usepackage{ulem}
\usepackage{nicefrac}
\usepackage{tikz}
\newtheorem{theorem}{Theorem}[section]
\newtheorem{lemma}[theorem]{Lemma}
\newtheorem{proposition}[theorem]{Proposition}
\newtheorem{corollary}[theorem]{Corollary}
\newtheorem{definition}[theorem]{Definition}
\newtheorem{example}[theorem]{Example}
\newtheorem{remark}[theorem]{Remark}
\newtheorem{question}{Question}
\newtheorem{examples}[theorem]{Examples}
\newtheorem{remarks}[theorem]{Remarks}
\newtheorem{claim}{Claim}
\renewcommand*{\theclaim}{\Alph{claim}}
\newcommand{\evid}[1]{\textsf{#1}}
\providecommand\llb{\llbracket}
\providecommand\rrb{\rrbracket}
\providecommand\spa{{\rm Span}}
\providecommand\I{{\mathcal{I}}}
\providecommand\Mon{{\mathcal{M}{\rm on}}}
\DeclareMathOperator{\mdeg}{\mathsf {m-deg}}
\DeclareMathOperator{\ini}{\mathsf{in}}
\begin{document}
\title[On atoms and their density in the monoid of monomial ideals]{On atoms and their density in the monoid of monomial ideals}
\author{Nikola Bogdanovic; Laura Cossu}
\date{}
\maketitle
\noindent\emph{Source and licence.} On atoms and their density in the monoid of monomial ideals. Version arXiv:2609.17202v1 (CC BY 4.0). This material is distributed under the Creative Commons Attribution 4.0 International licence, http://creativecommons.org/licenses/by/4.0/, as recorded in the arXiv OAI licence metadata for this exact version and shown on the abstract page https://arxiv.org/abs/2609.17202v1. Full rights notice: licenses/arxiv-2609.17202-CC-BY-4.0.txt.\par\smallskip
\noindent\textit{Editorial scope.} Counted unit: the original \texttt{theorem} environment \texttt{thm: 4-generated} at source lines 337--349 together with its complete proof at lines 351--469; the delivered document keeps source lines 132--470 so that every referenced label and every citation resolves inside it. Native editable source is the paper's own concatenated TeX; the AMS wrapper is editorial formatting only. No publisher class, upstream script or unrelated artwork is redistributed.
\begin{lemma}\label{lem: mon_gen}
Let $I\in\Mon(R)$, and $\mathcal G$ be a set of monomials in $I$. Then $\mathcal G$ generates $I$ if and only if for each monomial $v\in I$ there exists $u\in\mathcal G$ such that $u\mid v$.
\end{lemma}

We will also need the following elementary multiplicative property of the greatest common divisor of the minimal generators of a monomial ideal.
\begin{lemma}\label{lem: gcd IJ}
    For any $I,J\in\Mon(R)$, we have
    \[ \gcd(G(IJ)) = \gcd(G(I))\cdot\gcd(G(J)). \]
\end{lemma}
\begin{proof}
    Write $I=\gcd(G(I))\cdot I'$ and $J=\gcd(G(J))\cdot J'$, where $I',J'\in\Mon(R)$ and $\gcd(G(I'))=\gcd(G(J'))=1$. Clearly, we have $IJ = \gcd(G(I))\cdot \gcd(G(J))\cdot I'J'$ and $\gcd(G(IJ)) = \gcd(G(I))\cdot \gcd(G(J))\cdot \gcd(G(I'J'))$, so to conclude it suffices to show that $\gcd(G(I'J'))=1$. To achieve this, we prove that, for all $k\in [1,N]$, there is a minimal generator of $I'J'$ which is not divisible by $X_k$. Indeed, as $\gcd(G(I'))=1$ and $\gcd(G(J'))=1$, there are $f_k\in G(I')$ and $g_k\in G(J')$ such that $X_k$ divides neither $f_k$ nor $g_k$. As $f_kg_k\in I'J'$, there exists $h_k\in G(I'J')$ dividing $f_kg_k$, whence $X_k$ does not divide $h_k$.
\end{proof}

For a monomial ideal $I\in \Mon(R)$, we define its {\it min-degree} by:
\[\mdeg{I}=\min\{\deg(u) \; \colon \; u\in G(I)\}.\]
The min-degree is additive with respect to ideal multiplication: for all $I,J\in\Mon(R)$,
\[\mdeg(IJ)=\mdeg(I)+\mdeg(J).\]
Indeed, the ideal $IJ$ is generated by the monomials $uv$, with $u\in G(I)$ and $v\in G(J)$. Hence every element $w\in G(IJ)$ is divisible by some such product $uv$, and therefore $\deg(w)\geq \deg(u)+\deg(v)\geq \mdeg(I)+\mdeg(J)$. Thus,\[\mdeg(IJ)\geq \mdeg(I)+\mdeg(J).\]
Conversely, let $u\in G(I)$ and $v\in G(J)$ have degrees $\mdeg(I)$ and $\mdeg(J)$, respectively. Since $uv\in IJ$, some $w\in G(IJ)$ divides $uv$, and so
\[\mdeg(IJ)\leq \deg(w)\leq \deg(uv)=\mdeg(I)+\mdeg(J).\]

\section{Factorizations, lengths, and atoms in $\Mon(R)$}

Let $R=K[X_1,\ldots, X_N]$, with $K$ a field and $N\ge 2$. In this section, we study atoms of $\Mon(R)$ with a small number of minimal generators. More precisely, we give a complete characterization of atoms $I\in\Mon(R)$ with $\mu(I)\le 4$.

Before turning to this classification, we establish some basic factorization properties of principal monomial ideals and derive from them a realization result for sets of lengths. 

Recall that a \textit{factorization} of $I\in\Mon(R)$ is a formal product $z=A_1\cdots A_k$, such that \[I=A_1\cdots A_k,\]
where $A_1,\ldots,A_k$ are atoms of $\Mon(R)$, that is, proper monomial ideals of $R$ that cannot be written as a product of two proper monomial ideals. The integer $k$ is called the \textit{length} of $z$.
We denote by $\mathsf Z_{\Mon(R)}(I)$ the set of factorizations of $I$, where factorizations differing only by a permutation of the factors are identified. We also allow the empty factorization, which is the unique factorization of $R$ and has length $0$. The \textit{set of lengths} of $I$ is then defined by
\[
\mathsf L_{\Mon(R)}(I):=\{k\in\mathbb N : I \text{ admits a factorization of length } k\}.
\] 
It was proved in \cite[Theorem 4.7]{BCK} that $\Mon(R)$ is a $\mathsf{BF}$-monoid; equivalently, every element of $\Mon(R)$ admits a factorization and its set of lengths is finite.

Finally, recall that a monomial ideal $I$ is called a \textit{strong atom} (or an \textit{absolutely irreducible}) of $\Mon(R)$ if it is an atom and
\[\mathsf Z_{\Mon(R)}(I^n)=\{I^n\} \qquad\text{for every }n\ge 1.\]


\begin{lemma}\label{lem: princ div closed}
    Let $\underline X^a\in \mathcal M(R)$ with $a>\underline 0$ and suppose $\langle \underline X^a\rangle = JK$ for some $J,K\in\Mon(R)\setminus\{R\}$. Then both $J$ and $K$ are generated by a single non-constant monomial. In other words, the submonoid of $\Mon(R)$ consisting of principal monomial ideals is divisor-closed. In particular:
    \begin{enumerate}[label=\textup{(\arabic{*})}, mode=unboxed]
    \item\label{lem: princ div closed 1} $\mathsf Z_{\Mon(R)}(\langle \underline X^a\rangle) = \{ \langle X_1\rangle^{a(1)}\dots\langle X_N\rangle^{a(N)}\}$;
    \item\label{lem: princ div closed 2} for any $i\in [1,N]$, the ideal $\langle X_i\rangle$ is a strong atom of $\Mon(R)$.
    \end{enumerate}
\end{lemma}

\begin{proof}
    Let $G(J) = \{\underline X^{\beta_1}, \dots, \underline X^{\beta_{\mu(J)}}\}$, $G(K) = \{ \underline X^{\gamma_1}, \dots, \underline X^{\gamma_{\mu(K)}}\}$. We claim that $\mu(J)=\mu(K)=1$. Without loss of generality, we can assume $a=\beta_1+\gamma_1$. Suppose $\mu(K) \ge 2$: since $\underline X^{\beta_1+\gamma_2}\in \langle \underline X^a\rangle$, necessarily $a\le \beta_1+\gamma_2$. Combined with the previous equality, we obtain $\gamma_1\le \gamma_2$, implying $\underline X^{\gamma_1} \mid \underline X^{\gamma_2}$, contradicting the minimality of $G(K)$. Hence $\mu(K)=1$. By symmetry, $\mu(J)=1$ as well. Moreover, since $J,K\ne R$, their generators are nonconstant. This proves that the submonoid of $\Mon(R)$ consisting of principal monomial ideals is divisor-closed. Statement \ref{lem: princ div closed 1} follows from the fact that the monoid of monomials $\mathcal M(R)$ is free abelian on $X_1, \ldots,X_N$. Let $i\in [1,N]$ and take $a\in\mathbb N^N$ such that $a(i)=n$ and $a(j)=0$ for all $j\in [1,N]\setminus \{i\}$. It then follows from \ref{lem: princ div closed 1} that $\mathsf Z_{\Mon(R)}(\langle X_i^n\rangle) = \{ \langle X_i\rangle^{n}\}$, i.e., that $\langle X_i\rangle$ is a strong atom of $\Mon(R)$. This proves $\ref{lem: princ div closed 2}$.
    \end{proof}

\begin{proposition}\label{prop: princ factors}
    Let $I\in\Mon(R)$ and $\underline X^a\in \mathcal M(R)$. Then
    \[ \mathsf Z_{\Mon(R)} (\langle \underline X^a\rangle I) = \mathsf Z_{\Mon(R)}(\langle \underline X^a\rangle) \mathsf Z_{\Mon(R)}(I) = \langle X_1\rangle^{a(1)}\dots\langle X_N\rangle^{a(N)} \mathsf Z_{\Mon(R)}(I)\]
\end{proposition}
\begin{proof}
    If $a=\underline 0$, then $\langle \underline X^a\rangle=R$ and the claim is trivial. Assume then $a>\underline 0$. By Lemma \ref{lem: princ div closed}\ref{lem: princ div closed 1}, it suffices to prove the first equality. We do it for $\underline X^a = T\in\{X_1,\dots,X_N\}$, as the general case follows by iteration. It is clear that adjoining $\langle T\rangle$ to any factorization of $I$ yields a factorization of $\langle T\rangle I$, so it remains to prove the converse.

    Suppose that 
    \[\langle T\rangle I = U_1\cdots U_\ell,\]
    where $U_1,\dots,U_\ell$ are atoms of $\Mon(R)$. We claim that there exists a $j \in [1,\ell]$ such that $T\mid v$ for every $v\in G(U_j)$. Otherwise, for every $j\in[1,\ell]$ we could choose $v_j\in G(U_j)$ such that $T\nmid v_j$. Therefore,
    \[v_1\cdots v_\ell\in  U_1\cdots U_\ell=\langle T\rangle I\]
    while $T\nmid v_1\cdots v_\ell$, contradicting the fact that every monomial in $\langle T\rangle I$ is divisible by $T$. 

    Without loss of generality, assume that $T\mid v$ for every $v\in G(U_1)$. Then $U_1=\langle T\rangle U_1'$ for some $U_1'\in\Mon(R)$. Since $U_1$ is an atom, and $\langle T\rangle\ne R$, we must have $U_1'=R$, and so $U_1=\langle T\rangle$. Therefore,
    \[\langle T\rangle I = \langle T\rangle U_2\cdots U_\ell.\]
    As nonzero principal ideals are cancellative elements in $\mathcal I(R)$ (and so in $\Mon(R)$) by \cite{An-Ro97}, we obtain \[I = U_2\cdots U_\ell.\] Thus, every factorization of $\langle T\rangle I$ is obtained by adjoining $\langle T\rangle$ to a factorization of $I$, and this proves the claim.
\end{proof}

The following corollary is an immediate consequence of the above proposition.
\begin{corollary}
    For any $I\in\Mon(R)$, let $d=\gcd{(G(I))}$ and $J=\frac {1}{d} I$. Then, $J\in\Mon(R)$ and we have
    \[ \mathsf Z_{\Mon(R)}(I) = \mathsf Z_{\Mon(R)}(\langle d\rangle) \mathsf Z_{\Mon(R)}(J).  \]
\end{corollary}

Proposition \ref{prop: princ factors} also yields a translation property for sets of lengths: multiplying by a principal monomial ideal shifts all factorization lengths by the length of its unique factorization. This gives the following realization result.

\begin{proposition}\label{prop: relization of SOL}
  Let $m,n\in \mathbb N$ such that $2\le m\le n$. Then, there exists a monomial ideal $I\in \Mon(R)$ such that $\mathsf L_{\Mon(R)}(I)=[m,n]$.
\end{proposition}
\begin{proof}
It was proved in \cite[Theorem 4.7(3)]{BCK} that, for every integer $k\ge 2$,
\[ \mathsf L_{\Mon(R)}(\langle X_1,X_2\rangle^k)=[2,k].\]
Then, given \(2\le m\le n\), set $k=n-m+2$ and multiply $\langle X_1,X_2\rangle^k$ by $\langle X_1\rangle^{m-2}$. As $\langle X_1\rangle^{m-2}$ has a unique factorization of length $m-2$, we get from Proposition \ref{prop: princ factors} that
\[ \mathsf L_{\Mon(R)} \bigl(\langle X_1\rangle^{m-2}\langle X_1,X_2\rangle^{n-m+2}\bigr) = (m-2)+[2,n-m+2] = [m,n].\]
\end{proof}

\subsection{Characterizing \lq\lq small atoms\rq\rq}\label{char atom mu le 4}
In this section we provide a complete characterization of the atoms of $\Mon(R)$ with at most four minimal generators. The only atoms among the principal monomial ideals of $R$ are $\langle X_1\rangle, \dots, \langle X_N\rangle$ by Lemma \ref{lem: princ div closed}. The following results give necessary and sufficient conditions for $I=\langle\underline X^{a_1},\dots,\underline X^{a_{\mu(I)}}\rangle$ to be a (non-)atom in $\Mon(R)$, when $2\le \mu(I)\le 4$. We start with a couple of remarks establishing the theoretical and notational framework for our discussion, which will be used in the subsequent results without further comment.

\begin{remark}\label{rem: E1 E2}
    Let $I,J,K\in\Mon(R)\setminus\{R\}$ with $G(I) = \{\underline X^{a_1},\dots,\underline X^{a_{\mu(I)}}\}$, $G(J) = \{ \underline X^{\beta_1},\dots,\underline X^{\beta_{\mu(J)}} \}$, and $G(K) = \{ \underline X^{\gamma_1},\dots,\underline X^{\gamma_{\mu(K)}} \}$. By the minimality of the generating sets $G(I),G(J),G(K)$, we have that $a_i\perp a_j$ for every $i\neq j$, $\beta_k\perp\beta_\ell$ for $k\neq \ell$, and $\gamma_r\perp\gamma_s$ for $r\neq s$.

    With this notation, by Lemma \ref{lem: mon_gen}, the equality $I=JK$ holds if and only if the following two conditions are satisfied:
    \begin{enumerate}[label=\textup{\bf(E\arabic{*})}, mode=unboxed]
        \item\label{eq: I in JK} for every $i\in [1,\mu(I)]$, there exist $j_i\in [1,\mu(J)]$ and $k_i\in [1,\mu(K)]$ such that $a_i = \beta_{j_i}+\gamma_{k_i}$;
        \item\label{eq: JK in I} for every $j\in [1,\mu(J)]$ and every $k\in [1,\mu(K)]$, there is an $i\in [1,\mu(I)]$ such that $a_i \le \beta_j+\gamma_k$.
    \end{enumerate}
\end{remark}

\begin{remark}\label{rem: mon1 R}
    Let $I$ be a monomial ideal with $G(I) = \{\underline X^{a_1},\dots,\underline X^{a_{\mu(I)}}\}$. When studying conditions for $I$ to be an atom, we may assume that $I\neq R$, since atoms must be non-units. Moreover, in view of the preceding discussion, we may assume that $\gcd(G(I))=1$; otherwise, $I$ is divisible by a nontrivial principal monomial ideal. In terms of the exponent vectors, these conditions are equivalent to assuming $a_i> \underline 0$ for every $i\in [1,\mu(I)]$ and $\bigwedge_{i=1}^{\mu(I)}a_i=\underline 0.$ We denote by
    \[
        \Mon_1(R):=
        \{I\in\Mon(R):\gcd(G(I))=1\}
    \]
    the submonoid of $\Mon(R)$ consisting of monomial ideals with relatively prime minimal generators. By Lemma~\ref{lem: gcd IJ}, it is divisor-closed: indeed, if $IJ\in\Mon_1(R)$ for some $I,J\in\Mon(R)$, then
    \[
        1=\gcd(G(IJ))=\gcd(G(I))\gcd(G(J)),
    \]
    and hence
    \[
        \gcd(G(I))=\gcd(G(J))=1.
    \]
    Therefore, if $I=JK$ and $I\in\Mon_1(R)$, then, keeping the notation of Remark \ref{rem: E1 E2}, $\bigwedge_{j=1}^{\mu(J)} \beta_j = \bigwedge_{k=1}^{\mu(K)} \gamma_k = \underline 0$.
    Finally, $R$ is the unique principal ideal belonging to $\Mon_1(R)$.
    Indeed, if $\mu(I)=1$, say $G(I)=\{\underline X^{a}\}$, then $\gcd(G(I))=1$ if and only if $a=0$, equivalently, if and only if $I=R$. Consequently,
    \[
        I\in\Mon_1(R)\setminus\{R\}
        \quad\Longrightarrow\quad
        \mu(I)\ge 2.
    \]
\end{remark}

\bigskip


When $\mu(I)=2$, the next theorem shows that $I$ is an atom if and only if $I\in \Mon_1(R)$:

\begin{theorem}\label{thm: 2-gen}
    Let $I\in \Mon(R)$ with $G(I) = \{\underline X^{a_1},\underline X^{a_{2}}\}$. Then $I$ is not an atom of $\Mon(R)$ if and only if $a_1\wedge a_2>\underline 0$.
\end{theorem}
\begin{proof}
    Since $a_1,a_2\in\mathbb N^N$, we have $a_1\wedge a_2\in\mathbb N^N$, i.e., $a_1\wedge a_2\ge \underline 0$. If $a_1\wedge a_2 > \underline 0$, then, using that $a_1\perp a_2$, it is easy to see that
   \[J=\langle \underline X^{a_1\wedge a_2}\rangle \quad \text{ and }\quad K=\langle \underline X^{a_1-a_1\wedge a_2}, \underline X^{a_2-a_1\wedge a_2} \rangle \]
   are non-units of $\Mon(R)$. Thus, $I=JK$ is not an atom.
    On the other hand, assume that $I=JK$ for some $J,K\in\Mon(R)\setminus\{R\}$, and keep the notation of Remark \ref{rem: E1 E2}. By condition \ref{eq: I in JK}, there are $j_1,j_2\in [1,\mu(J)]$ and $k_1,k_2\in [1,\mu(K)]$ such that $a_1 = \beta_{j_1}+\gamma_{k_1}$ and $a_2=\beta_{j_2}+\gamma_{k_2}$. By condition \ref{eq: JK in I}, we have $a_i\le \beta_{j_1}+\gamma_{k_2}$ for either $i=1$ or $i=2$. Since $a_i=\beta_{j_i}+\gamma_{k_i}$, we obtain
    \[ \beta_{j_i}+\gamma_{k_i} \le \beta_{j_1}+\gamma_{k_2}, \]
    whence either $\gamma_{k_1}\le \gamma_{k_2}$ or $\beta_{j_2}\le \beta_{j_1}$. By minimality of $G(J)$ and $G(K)$, we obtain either $k_1=k_2$ or $j_2=j_1$. In both cases, by condition \ref{eq: I in JK} we have $a_1\wedge a_2>\underline 0$.
\end{proof}


The next theorem gives a complete characterization of (non-)atoms among the elements of $\Mon(R)$ with $3$ minimal generators. 
\begin{theorem}\label{thm: 3-generated}
Let $I\in \Mon_1(R)$ with $G(I) = \{\underline X^{a_1},\underline X^{a_{2}}, \underline X^{a_3}\}$. Then $I$ is not an atom of $\Mon(R)$ if and only if there exists a permutation $\sigma\in S_3$ such that $a_{\sigma(1)}+a_{\sigma(2)}\ge 2a_{\sigma(3)}$.
\end{theorem}
\begin{proof}
Assume that $I$ is not an atom of $\Mon(R)$. As $I\ne R$, this means that $I=JK$ for some $J,K\in\Mon(R)\setminus\{R\}$. Using the notation of Remark \ref{rem: E1 E2}, 
by \ref{eq: I in JK}, there exist $j_1,j_2,j_3\in [1,\mu(J)]$ and $k_1,k_2,k_3\in [1,\mu(K)]$ such that $a_i=\beta_{j_i}+\gamma_{k_i}$ for every $i\in [1,3]$. We proceed by analyzing the cardinality of the index sets $\{j_1,j_2,j_3\}$ and $\{k_1,k_2,k_3\}$.

\medskip
\noindent\textbf{Case 1: $\lvert\{j_1,j_2,j_3\}\rvert = 1$.} 
This is impossible, as $j_1=j_2=j_3$ implies $\underline{0} < \beta_{j_1} \le a_1 \wedge a_2 \wedge a_3$ by \ref{eq: I in JK}, a contradiction.

\medskip
\noindent\textbf{Case 2: $\lvert\{j_1,j_2,j_3\}\rvert = 2$.} 
Assume without loss of generality that $j_1 \neq j_2 = j_3$, which implies $k_2 \neq k_3$ (otherwise $a_2=a_3$ by \ref{eq: I in JK}). 

\begin{itemize}
    \item Assume $k_1=k_2$. Condition \ref{eq: JK in I} forces $a_i \le \beta_{j_1} + \gamma_{k_3}$ for some $i\in[1,3]$. 
    By \ref{eq: I in JK}, both $i=1$ and $i=3$ yield immediate contradictions ($\gamma_{k_1} \le \gamma_{k_3}$ and $\beta_{j_3} \le \beta_{j_1}$, respectively). Thus $a_2 \le \beta_{j_1} + \gamma_{k_3}$, which leads to:
    \[ a_1+a_3 = \beta_{j_1}+\gamma_{k_1}+\beta_{j_3}+\gamma_{k_3}=\beta_{j_1}+\gamma_{k_2}+\beta_{j_2}+\gamma_{k_3} \ge 2a_2 .\]
    \item The subcase $k_1 = k_3$ is perfectly analogous to the previous one.
    \item Assume that $k_1,k_2, k_3$ are pairwise distinct. By \ref{eq: JK in I}, we have $a_i\le \beta_{j_2}+\gamma_{k_1}$ for some $i\in [1,3]$. If $i=1$, we obtain $a_1=\beta_{j_1}+\gamma_{k_1}\le \beta_{j_2}+\gamma_{k_1}$ by \ref{eq: I in JK}, implying $\beta_{j_1}\le \beta_{j_2}$, against $j_1\neq j_2$. Similarly, one finds a contradiction in each of the subcases $i=2$ and $i=3$.
\end{itemize}

\medskip
\noindent\textbf{Case 3:$\lvert\{ j_1,j_2,j_3\}\rvert=3$.}
The subcase $\lvert\{k_1,k_2,k_3\}\rvert \le 2$ is ruled out by symmetric arguments to Cases 1 and 2 (exchanging the roles of $J$ and $K$). Assume then $\lvert\{k_1,k_2,k_3\}\rvert = 3$. By \ref{eq: JK in I}, for any pair $(j_r, k_s)$ with $r \neq s$, there exists $a_i \le \beta_{j_r} + \gamma_{k_s}$. If $i=r$ or $i=s$, we obtain a contradiction via \ref{eq: I in JK} (e.g., $a_1 \le \beta_{j_1}+\gamma_{k_2} \implies \gamma_{k_1} \le \gamma_{k_2}$). 
Thus, the index $i$ must be the unique element in $\{1,2,3\} \setminus \{r,s\}$. By cyclic permutations, this yields:
\begin{equation}\label{eq: mu3 system}
    a_3 \le \beta_{j_1}+\gamma_{k_2}, \quad a_1 \le \beta_{j_2}+\gamma_{k_3}, \quad a_2 \le \beta_{j_3}+\gamma_{k_1}.
\end{equation}
but also
\begin{equation}\label{eq: mu3 system 2}
    a_3 \le \beta_{j_2}+\gamma_{k_1}, \quad a_1 \le \beta_{j_3}+\gamma_{k_2}, \quad a_2 \le \beta_{j_1}+\gamma_{k_3}.
\end{equation}
Summing the relations \eqref{eq: mu3 system} gives:
\[
    a_1+a_2+a_3 \le (\beta_{j_1}+\gamma_{k_2}) + (\beta_{j_2}+\gamma_{k_3}) + (\beta_{j_3}+\gamma_{k_1}) = a_1+a_2+a_3,
\]
which forces all inequalities in \eqref{eq: mu3 system} to hold as equalities. In particular, we have $a_3 = \beta_{j_1}+\gamma_{k_2}$. Proceeding analogously with the relations in \eqref{eq: mu3 system 2}, we get $a_2 = \beta_{j_1}+\gamma_{k_3}$. Thus, combining the two above equalities with $a_1 = \beta_{j_1}+\gamma_{k_1}$ from \ref{eq: I in JK}, we obtain $\underline{0} < \beta_{j_1} \le a_1 \wedge a_2 \wedge a_3$, a contradiction.

Vice versa, let us show that $a_{\sigma(1)}+a_{\sigma(2)}\ge 2a_{\sigma(3)}$ for some $\sigma\in S_3$ is also a sufficient condition for $I=\langle \underline X^{a_1}, \underline X^{a_2}, \underline X^{a_3} \rangle$ to be a non-atom of $\Mon(R)$. Up to renaming the exponents of the generators of $I$, we can assume that $a_1+a_2\ge 2a_3$. Note that $a_1\wedge a_3>\underline 0$, otherwise $a_2\ge 2a_3 > a_3$, against $a_2\perp a_3$. Analogously, $a_2\wedge a_3>\underline 0$. We can therefore define the proper ideals:
\[
    J = \langle \underline{X}^{a_1 \wedge a_3}, \underline{X}^{a_2-a_3+a_1 \wedge a_3} \rangle, \qquad
    K = \langle \underline{X}^{a_1 - a_1 \wedge a_3}, \underline{X}^{a_3 - a_1 \wedge a_3} \rangle.
\]
Indeed, no generator exponent equals $\underline{0}$ because $a_1, a_3 > a_1 \wedge a_3$ (otherwise $a_3\ge a_1$ or $a_1\ge a_3$), and for every component $i \in [1,N]$:
\[
    a_2(i) + (a_1 \wedge a_3)(i) = a_2(i)+a_1(i)\ge 2a_3(i) \ge a_3(i),
\]
or
\[
    a_2(i) + (a_1 \wedge a_3)(i) = a_2(i)+a_3(i) \ge a_3(i),
\]
with at least one strict inequality since $a_2 \wedge a_3 > \underline{0}$. Finally, we verify that $JK = I$. Clearly $\underline{X}^{a_1}, \underline{X}^{a_2}, \underline{X}^{a_3} \in JK$. The remaining product generator satisfies:
\[
    (a_2 + a_1 \wedge a_3 - a_3) + (a_1 - a_1 \wedge a_3) = a_2 + a_1 - a_3 \ge 2a_3 - a_3 = a_3,
\]
hence $\underline{X}^{a_2+a_1-a_3} \in \langle \underline{X}^{a_3} \rangle \subseteq I$, completing the proof.
\end{proof}

We conclude this section by characterizing the non-atoms of $\Mon(R)$ with $4$ minimal generators. The proof proceeds by a case analysis and relies on arguments analogous to those detailed in the proof of the previous theorem. To avoid unnecessary repetition, we omit routine symmetric arguments and give details only when new features arise.


\begin{theorem}\label{thm: 4-generated}
Let $I\in \Mon_1(R)$ with $G(I) = \{\underline X^{a_1},\underline X^{a_{2}}, \underline X^{a_3}, \underline X^{a_4}\}$. Then $I$ is not an atom of $\Mon(R)$ if and only if it satisfies one of the following conditions, up to a relabeling of the exponents of its generators:
\begin{enumerate}[label=\textup{(c\arabic{*})}, mode=unboxed]
 \item \label{mu4 c1} $a_1+a_3=a_2+a_4$; 
 \item \label{mu4 c2} $a_1+a_3\ge 2a_4$ and $a_2+a_4\ge 2a_1$; 
 \item \label{mu4 c3} $a_1\wedge a_2\wedge a_3>\underline 0$ and one of the following holds:
        \begin{enumerate}[label=\textup{(\roman{*})}, mode=unboxed]
            \item \label{mu4 c31} $a_2+a_4\ge a_1+a_3$ and $a_3+a_4\ge 2a_1$;
            \item \label{mu4 c32} $a_2+a_4\ge 2a_1$ and $a_3+a_4\ge a_1+a_2$;
            \item \label{mu4 c33} $a_2+a_4\ge 2a_1$ and $a_3+a_4\ge 2a_1$.
        \end{enumerate}
    \end{enumerate}
\end{theorem}

\begin{proof}
Suppose that $I$ is not an atom of $\Mon(R)$, i.e., $I=JK$ for some $J,K\in\Mon(R)\setminus\{R\}$. We use the notation introduced in Remark \ref{rem: E1 E2}: for every $i\in[1,4]$ we choose indices $j_i,k_i$ such that
\[
a_i=\beta_{j_i}+\gamma_{k_i}.
\]
Recall that whenever $j_r\neq j_s$, condition \ref{eq: JK in I} applied to the crossed sum $\beta_{j_r}+\gamma_{k_s}$ cannot hold for $a_s$, as this would contradict the minimality of $G(J)$. Similarly, whenever $k_r\neq k_s$, this sum cannot dominate $a_r$ due to the minimality of $G(K)$. Consequently, if $j_r\neq j_s$ and $k_r\neq k_s$, the crossed sum $\beta_{j_r}+\gamma_{k_s}$ must dominate one of the remaining generators. We will use this observation repeatedly without further mention.

Up to relabeling the $a_i$'s and the generators of $J$ and $K$, and possibly interchanging $J$ and $K$, we may assume $|\{j_1,j_2,j_3,j_4\}|\le|\{k_1,k_2,k_3,k_4\}|$. We then distinguish the following cases, depending on the cardinality and multiplicity pattern of the index set $\{j_i\}_{i=1}^4$.
\begin{enumerate}[label=\textup{(\arabic{*})}, mode=unboxed]
        \item\label{mu4 4000} $j_1 = j_2 = j_3 = j_4$;
        \item\label{mu4 3100}  $j_1=j_2=j_3\neq j_4$;
        
        \item\label{mu4 2200} $j_1=j_2\ne j_3=j_4$;
        
        \item\label{mu4 2110} $\lvert\{j_1,j_2,j_3,j_4\}\rvert = 3,\ j_1=j_2$;
        
        \item\label{mu4 1111} $\lvert\{j_1,j_2,j_3,j_4\}\rvert=4$.
\end{enumerate}
Recall that, by the minimality of $G(I)$, if $j_r=j_{s}$ for distinct $r,s\in[1,4]$, then $k_r\ne k_{s}$.

\smallskip
\textbf{Case \ref{mu4 4000}.} This case cannot occur, as it implies $a_1\wedge a_2\wedge a_3\wedge a_4 \ge \beta_{j_1} > \underline 0$, against our assumptions.

\smallskip
\textbf{Case \ref{mu4 3100}.} Suppose $j_1=j_2=j_3\neq j_4$. If $k_4\notin\{k_1,k_2,k_3\}$, then the crossed sum $\beta_{j_1}+\gamma_{k_4}$ cannot dominate any $a_i$, contradicting \ref{eq: JK in I}. Hence, after relabeling, we may assume $k_4=k_1$. Then $a_1\wedge a_2\wedge a_3\ge \beta_{j_1}>\underline 0.$ By \ref{eq: JK in I} there are $i\in \{1,3\}$ and $\ell\in\{1,2\}$ such that
\[\beta_{j_4}+\gamma_{k_2}\ge a_i \quad\text{ and }\quad \beta_{j_4}+\gamma_{k_3}\ge a_\ell.\] The choices $(i, \ell) = (3, 1), (1, 2), (1, 1)$ yield, respectively, \ref{mu4 c3}\ref{mu4 c31}, \ref{mu4 c3}\ref{mu4 c32}, and \ref{mu4 c3}\ref{mu4 c33}. The choice $(i,\ell)=(3,2)$ yields $a_4\ge a_1$, a contradiction.

\smallskip
\textbf{Case \ref{mu4 2200}.} Suppose that $j_1=j_2\neq j_3=j_4$. If the four $k_i$'s are distinct, the forced crossed inequalities imply simultaneously
\[
a_2+a_3\ge a_1+a_4,\qquad
a_2+a_4\ge a_1+a_3,
\]
and hence $a_2\ge a_1$, a contradiction. If, after relabeling, $k_1=k_4$ and $k_2\neq k_3$, then the crossed sums give
\[
a_1+a_3\ge 2a_4,\qquad a_2+a_4\ge 2a_1,
\]
which is condition \ref{mu4 c2}. Finally, if $k_1=k_4$ and $k_2=k_3$, then immediately
\[
a_1+a_3=a_2+a_4,
\]
which is condition \ref{mu4 c1}.

\smallskip
\textbf{Case \ref{mu4 2110}.} Assume that $j_1=j_2$ and $j_3,j_4$ are distinct from each other and from $j_1$. If the four $k_i$'s are distinct, a repeated application of the crossed-sum observation forces a chain of inequalities in which every alternative except one leads to a forbidden comparison between two of the $a_i$'s. In the remaining alternative, the inequalities are forced into equalities, yielding
\[
a_1\wedge a_2\wedge a_3\wedge a_4\ge \beta_{j_1}>\underline 0,
\]
a contradiction. If, on the other hand, the set $\{k_i\}_{i=1}^4$ has cardinality $3$, say $k_3=k_4\notin\{k_1,k_2\}$, then the crossed sum $\beta_{j_1}+\gamma_{k_3}$ cannot dominate any of the $a_i$, again a contradiction. In the remaining subcase, after relabeling, $k_2=k_3$ and $k_4\notin\{k_1,k_2\}$. In this case the crossed sums first give
\[
a_2+a_4\ge a_1+a_3.
\]
The remaining relevant crossed sum $\beta_{j_3} +\gamma_{k_1}$ cannot dominate $a_2$, otherwise $a_4\ge a_2$, thus it must dominate $a_4$ and the previous inequality becomes an equality:
\[
a_2+a_4=a_1+a_3.
\]
Thus condition \ref{mu4 c1} holds.
%A more careful inspection even shows that this subcase is impossible, but this is not needed.

\textbf{Case \ref{mu4 1111}.} In this case, all the $j_i$'s are pairwise distinct, and hence so are all the $k_i$'s. For $i\neq \ell$, choose
\[
\nu(i,\ell)\in[1,4]\setminus\{i,\ell\}
\quad \text{ such that }\quad
\beta_{j_i}+\gamma_{k_\ell}\ge a_{\nu(i,\ell)}.
\]

We first claim that $\nu(i,\ell)=\nu(\ell,i)$
for all $i\neq \ell$. Indeed, if $\nu(i,\ell)=m$ and $\nu(\ell,i)=n$ are distinct, then $\{i,\ell,m,n\}=[1,4]$ and
\[
a_i+a_\ell
\ge a_m+a_n
\ge a_{\nu(m,n)}+a_{\nu(n,m)}.
\]
The possible choices for $\nu(m,n)$ and $\nu(n,m)$ either immediately contradict incomparability of the $a_i$'s, or force equality throughout. In the latter case one obtains identities of the form
\[
\beta_{j_i}+\gamma_{k_\ell}=a_m,\qquad
\beta_{j_i}+\gamma_{k_n}=a_\ell,\qquad
\beta_{j_i}+\gamma_{k_m}=a_n,
\]
and hence
\[
a_1\wedge a_2\wedge a_3\wedge a_4\ge \beta_{j_i}>\underline 0,
\]
again impossible. This proves the claim that $\nu(i,\ell)=\nu(\ell,i)$ for all $i\ne\ell$. It remains to consider the possible assignments of these values. Up to a relabeling of the indices, we may fix $\nu(1,2)$. A finite case analysis, repeatedly applying the crossed-sum observation, shows that every assignment either yields a forbidden comparison $a_r\ge a_s\qquad (r\ne s)$, or forces equality in a sum of crossed inequalities and consequently contradicts the minimality of $G(J)$ or $G(K)$. Therefore this case cannot occur. 

We have proved that if $I$ is not an atom, then one of \ref{mu4 c1}, \ref{mu4 c2}, \ref{mu4 c3} holds.

Conversely, suppose first that \ref{mu4 c1} holds, say $a_1+a_3=a_2+a_4.$ Set $d=a_1\wedge a_2$. Then
\[
I=
\langle \underline X^d, \underline X^{a_4-a_1+d}\rangle
\langle \underline X^{a_1-d},\underline X^{a_2-d}\rangle .
\]
Indeed, the four products give $a_1,a_2,a_4$, and $a_4-a_1+d+a_2-d=a_4-a_1+a_2=a_3.$ Moreover, the exponents involved are all $>\underline 0$: for instance, $d=\underline 0$ would imply $a_3\ge a_2$, and $a_4-a_1+d=\underline 0$ would imply $a_1\ge a_4$, both contradicting the minimality of $G(I)$.

Suppose next that \ref{mu4 c2} holds, say $a_1+a_3\ge 2a_4$ and $a_2+a_4\ge 2a_1.$ Set $d=a_1\wedge a_4$. Then
\[
I=
\langle \underline X^{a_1-d},\underline X^{a_4-d}\rangle
\langle \underline X^d,\underline X^{a_2-a_1+d},\underline X^{a_3-a_4+d}\rangle .
\]
The inequalities ensure that the extra products are absorbed:
\[
a_4+a_2-a_1\ge a_1,\qquad a_1+a_3-a_4\ge a_4.
\]
The same component-wise check as above shows that all displayed vector exponents are $>\underline 0$.

Finally, assume \ref{mu4 c3}. Put $d=a_1\wedge a_2\wedge a_3>\underline 0.$ Then
\[
I=
\langle \underline X^d,\underline X^{a_4-a_1+d}\rangle
\langle \underline X^{a_1-d},\underline X^{a_2-d},\underline X^{a_3-d}\rangle .
\]
The products with $\underline X^d$ give $a_1,a_2,a_3$, while the product $(a_4-a_1+d)+(a_1-d)$ gives $a_4$. The remaining two products have exponents
\[
a_4-a_1+a_2,\qquad a_4-a_1+a_3,
\]
and these dominate one of $a_1,a_2,a_3$ precisely by the alternatives listed in \ref{mu4 c3}. Thus they are absorbed by $I$. Again, the hypotheses and the incomparability of the $a_i$'s ensure that all exponents are $>\underline 0$. Hence in each case $I$ admits a nontrivial factorization, and so it is not an atom.
\end{proof}

\begin{thebibliography}{99}
\bibitem[An-Ro97]{An-Ro97} D.~D.~Anderson and M.~Roitman, \emph{A characterization of cancellation ideals}, Proc.~Am.~Math.~Soc. \textbf{125} (1997), No.~10, 2853--2854. %
\bibitem[BCK]{BCK} N. Bogdanovic, L. Cossu, M. A. Khadam, \emph{On the arithmetic of polynomial ideals}, J. Pure Appl. Algebra, to appear. Preprint version: arXiv:2510.24455. %
\end{thebibliography}
\end{document}
